\documentclass{article}
\usepackage[T1]{fontenc}

\begin{document}

ok i'll try to make as much sense as i can. ooh also first time using latex formally informally yay
\bigskip

kepler's third law states that the square of the orbital period of a planet (P) is directly proportional to the cube of the semi-major axis of its orbit (a). this means that if you know the distance of a planet from the sun, you can calculate how long it takes to complete one orbit around the sun.
\smallskip

mathematically,

\begin{equation}
P^2 \propto a^3
\end{equation}

when we compare two planets orbiting the same star/central body, it forms a ratio,

\begin{equation}
\frac{P_1^2}{a_1^3} = \frac{P_2^2}{a_2^3}
\end{equation}

if we measure the distances in AU, and the orbital periods in years, the constant of proportionality becomes 1 for all planets in our solar system,

\begin{equation}
P^2 = a^3
\end{equation}




\end{document}